\documentclass[10pt]{article}

\usepackage[utf8]{inputenc}

\usepackage[T1]{fontenc}

\usepackage{amsmath}

\usepackage{amsfonts}

\usepackage{amssymb}

\usepackage[version=4]{mhchem}

\usepackage{stmaryrd}

\usepackage{graphicx}

\usepackage[export]{adjustbox}

\graphicspath{ {./images/} }



\begin{document}

Если $f(m)$ - действительная аддитивная арифмметич. функция,



\[

A_{n}=\sum_{p \leqslant n} \frac{f(p)}{p}, B_{n}^{2}=\sum_{p}{ }^{\alpha} \frac{f^{2}\left(p^{\alpha}\right)}{p^{\alpha}},

\]



где суммы берутся по простым числам $p$ и по степеням простых чисел $p^{\alpha}$, то



\[

\frac{1}{n} \sum_{m=1}^{n}\left(f(m)-A_{n}\right)^{2} \leqslant B_{n}^{2}\left(\frac{3}{2}+\frac{c}{\ln n}\right),

\]



где $c$ - абсолютная константа. Следовательно, для любого $t>0$ и всех $m \leqslant n$, за исключением $<\left(\frac{3}{2}+\frac{c}{\ln n}\right) n t^{-2}$ чисел, имеет место неравенство



\[

\left|f(m)-A_{n}\right|<t B_{n}

\]



(аналог теоретико-вероятностного больиих чисел зако$н a)$. Для функции $\omega(m)$ это неравенство можно записать в виде



\[

|\omega(m)-\ln \ln n|<t \sqrt{\ln \ln n} .

\]



Пусть через $N_{n}(\ldots)$ обозначено число натуральных $m \leqslant n$, удовлетворяющих условиям, к-рые будут указываться в скобках вместо многоточия. Желая более точно охарактеризовать распределение значений действительных арифметич. функций $f(m)$, приходят к рассмотрению асимптотич. поведения частоты



\[

\frac{1}{n} N_{n}(f(m) \in E)

\]



при $n \rightarrow \infty$, где $E$ - любое борелевское множество. Среди асимптотич. законов для (3) наибольший интерес представляют законы двух типов: интегральные і локальвые.



Интегральные законы. Изучается асимптотич. поведение функции распределения



\[

F_{n}\left(C_{n}+D_{n} x\right)=\frac{1}{n} N_{n}\left(f(m)<C_{n}+D_{n} x\right),

\]



прпи $n \longrightarrow \infty$ и заданных $C_{n}, D_{n}$.



В случае арифометических аддитивных функциї ищутся условия, при к-рых $F_{n}\left(C_{n}+D_{n} x\right)$ стремится К нек-рой функции распределенія $F(x)$ во всех ее точках непрерывности. При этом, если $F(x)$ не вырождены, то $D_{n}$ обязательно должно стремиться к конечному (отличному от 0 ) или бесконечному пределу.



В случае конечного предела достаточно ограничиться рассмотрением $F_{n}\left(C_{n}+x\right)$. Для того чтобы $F_{n}\left(C_{n}+x\right)$ с какими-либо $C_{n}$ при $n \longrightarrow \infty$ имела невырожденное предельное распределение, необходимо и достаточно, чтобы $f(m)$ Імела вид



\[

f(m)=a \ln m+g(m)

\]



где $a$ - константа, а функция $g(m)$ удовлетворяет условиям



\[

\sum_{|g(p)| \geqslant 1} \frac{1}{p}<\infty, \sum_{|g(p)|<1} \frac{g^{2}(p)}{p}<\infty .

\]



При этом $C_{n}$ должны быть равными



\[

C_{n}=a \ln n+\sum_{\substack{p \leqslant n,|g(p)|<1}} \frac{g(p)}{p}+C+o(1),

\]



$C$-константа. Выбор $C_{n}$ однозначев с точностью до слагаемых $C+o(1)$. Предельно распределение является дискретным, когда $\sum_{f(p) \neq 0} \frac{1}{p}<\infty$, и непрерывным в противном случае.



В частности, $F_{n}(x)$ (случай $C_{n} \equiv 0$ ) тогда пl только тогда имеет предельное распределение, когда сходятся ряды



\[

\sum_{|f(p)| \geqslant 1} \frac{1}{p}, \sum_{|f(p)|<1} \frac{f(p)}{p}, \sum_{|f(p)|<1} \frac{f^{2}(p)}{p}

\]



(аналог теоретико-вероятностной теоремы о трех рядах).



Случай $D_{n} \longrightarrow \infty$ не исследован до конца. Ниже приведены нек-рые наиболее простые результаты, когда $C_{n}=A_{n}$ и $D_{n}=B_{n}$ определены формулами (2).



Если для всякого фиксированного $\varepsilon>0$



\begin{center}

\includegraphics[max width=\textwidth]{2023_07_09_c0184e6a865f9c443b70g-1}

\end{center}



при $n \rightarrow \infty$ (аналог условия Лпндеберга, см. Линдеберга - Феллера теорема), то



\[

F_{n}\left(A_{n}+B_{n} x\right) \rightarrow \frac{1}{\sqrt{2 \pi}} \int_{-\infty}^{x} e^{-u^{2} / 2} d u=\Phi(x)

\]



(нормальный закон). Если выполнено (4), то $B_{n}$ является медленно меняющейся функцией от $\ln n$ в смысле Карамата. Более того, если $B_{n}$ является такой функцией, то для справедливости (5) условие (4) является необходимым.



Пусть $B_{n}$ является медленно меняющейся функцией от $\ln n$. Для того чтобы $F_{n}\left(A_{n}+B_{n} x\right)$ сходилась к предельному распределению с дисперсией 1 , необходимо и достаточно, чтобы существовала такая неубывающая функция $V(u),-\infty<u<\infty, V(-\infty)=0, V(\infty)=1$, что при $n \rightarrow \infty$ для всех $u$, за исключением, быть может, $u=0$,



\[

\begin{aligned}

& B_{n}^{-2} \sum_{p^{\alpha} \leqslant n,} \frac{f^{2}\left(p^{\alpha}\right)}{p^{\alpha}} \rightarrow V(u) . \\

& f\left(p^{\alpha}\right)<u B_{n}

\end{aligned}

\]



Характеристич. функция $\varphi(t)$ предельного закона в случае его существования определяется формулой



\[

\varphi(t)=\exp \left(\int_{-\infty}^{\infty}\left(e^{i t u}-1-i t u\right)\right) u^{-2} d V(u) .

\]



Изучается быстрота сходимости к предельному закону. Так, напр., если $f(m)$ - сильно аддитивная функция II



\[

\mu_{n}=B_{n}^{-1} \max _{p \leqslant n}|f(p)| \rightarrow 0

\]



ири $n \rightarrow \infty$, то



\[

F_{n}\left(A_{n}+B_{n} x\right)=\Phi(x)+O\left(\mu_{n}\right)

\]



равномерно по $x$.



Для мультипликативных арифметич. функций имеют место аналогичныс результаты.



Јокальные законы. Изучается поведение при $n \rightarrow \infty$ частоты $\frac{1}{n} N_{n}(f(m)=c)$ при заданных $c$. В случае действительных аддитивных арифметит. функций әта частота всегда имеет предел, к-рый отличен от 0 липь для не более тем счетного множества значениї $c$. Пусть



\[

\lambda_{l}(f)=\lim _{n \rightarrow \infty} \frac{1}{n} N_{n}\left(f(m)=c_{l}\right)

\]



\begin{itemize}

  \item не равные нулю пределы, причем хотя бы один такой предел существует. Тогда

\end{itemize}



\[

\sum_{l} \lambda_{l}(f)=1

\]



II



\[

\sum_{l} \lambda_{l}(f) e^{i t c} l=\prod_{p}(1-1 / p) \sum_{\alpha=0}^{\infty} p^{-\alpha} e^{i t f\left(p^{\alpha}\right)} .

\]



Если $f(m)$ принимает лишь целые значения,



To



\[

\mu_{k}(f)=\lim _{n \rightarrow \infty} \frac{1}{n} N_{n}(/(m)=k)

\]



\[

\sum_{k} \mu_{k}(f)=1

\]





\end{document}